\documentclass[a4paper,11pt]{article}
\usepackage{exam-style-842}
\examinehead{842 信息化概论}{模拟试卷（一）参考答案}
\begin{document}

\answercover{842 信息化概论}{模拟试卷（一）参考答案}

\answerpart{一、名词解释（每题 5 分，共 30 分）}

\q \textbf{信息化}\\
信息化是充分利用信息技术、开发利用信息资源、促进信息交流和知识共享，从而提高经济增长质量、推动经济社会发展转型的历史进程。其内涵可从三方面把握：一是技术基础是信息技术，二是核心是信息资源的开发利用与共享，三是目标是实现由工业社会向信息社会的转变。信息化具有综合性、竞争性、渗透性、开放性；我国信息化的发展方向表现为知识型经济、网络化社会、数字化生活和服务型政府。数字化是信息化实现的必备条件，信息化是数字化的目标和归宿。
\begin{analysis}得分点：①“充分利用信息技术、开发利用信息资源、促进信息交流与知识共享”给 2 分；②“推动经济社会发展转型的历史进程”给 2 分；③补充“数字化是基础、信息化是目标”或“四性”给 1 分。常见失分：只答成“用计算机办公”，未点出信息资源开发利用这一核心；把信息化与数字化混为一谈而不作区分。\end{analysis}

\q \textbf{云计算}\\
云计算是通过网络按需向用户提供可配置计算资源（网络、服务器、存储、应用和服务）并实现按使用量付费的一种服务模式。其基本特征是按需自助服务、广泛的网络访问、资源池化、快速弹性伸缩和可计量的服务；关键技术是虚拟化技术、分布式存储与并行计算、多租户隔离与资源调度。按服务方式分为 IaaS（基础设施即服务）、PaaS（平台即服务）、SaaS（软件即服务），按部署方式分为公有云、私有云、混合云和社区云。
\begin{analysis}得分点：①按需提供、按量付费的服务模式（2 分）；②五大基本特征或资源池化、弹性（2 分）；③虚拟化等关键技术，或写出 IaaS/PaaS/SaaS 三层（1 分）。常见失分：把云计算等同于“远程服务器托管”或“大型机”，未体现虚拟化与弹性；把 SaaS 说成提供虚拟机。\end{analysis}

\q \textbf{大数据}\\
大数据是指规模巨大、类型多样、产生和处理速度快、价值密度低但总体价值高的数据集合，通常用 4V 特征描述，即数据量大（Volume）、类型多（Variety）、速度快（Velocity）、价值密度低（Value）。其技术体系包括数据采集、存储、计算分析与可视化四个层次；采集环节依赖传感器与网络日志，存储环节以分布式文件系统 HDFS 和 NoSQL 数据库为主，计算环节以 MapReduce、Spark 等并行计算框架为主，分析与可视化环节包括数据挖掘、统计分析和图表呈现。
\begin{analysis}得分点：①4V 特征答出 3 个以上（2 分）；②“价值密度低、总体价值高”的辩证表述（1 分）；③技术体系或 Hadoop、MapReduce、可视化等具体技术（2 分）。常见失分：只写“数据很多”，不答特征；把大数据与数据挖掘、云计算混作同一概念。\end{analysis}

\q \textbf{物联网}\\
物联网是通过射频识别、红外感应器、全球定位系统、激光扫描器等感知设备，按约定的协议把物品与互联网连接起来，进行信息交换和通信，以实现对物品的智能化识别、定位、跟踪、监控和管理的一种网络。其体系结构分为感知层、网络层和应用层；核心技术包括传感器与传感网技术、RFID 技术、网络通信技术和数据处理与智能控制技术。物联网的实质是“物物相联、信息互通、智能控制”，典型应用有农产品溯源、智能温室、智能家居和智能交通。
\begin{analysis}得分点：①感知设备＋网络协议＋智能化识别定位跟踪监控管理（3 分）；②三层体系结构（1 分）；③举出 RFID、传感器网络或农业溯源应用（1 分）。常见失分：只答成“传感器网络”，漏掉“物品与互联网互联”和“智能管理”的目的；三层结构顺序写错。\end{analysis}

\q \textbf{OSI 参考模型}\\
OSI 参考模型是国际标准化组织（ISO）提出的开放系统互连参考模型，它将网络通信功能划分为七层，自下而上依次为物理层、数据链路层、网络层、传输层、会话层、表示层和应用层。分层的目的是把复杂的网络通信问题分解为若干相对独立、接口清晰的子问题，实现模块化设计、易于标准化和互操作。各层的主要功能是：物理层传输比特流；数据链路层以帧为单位实现相邻结点间的可靠传输与介质访问控制；网络层负责分组转发与路由选择；传输层提供端到端的可靠或不可靠传输；会话层管理会话；表示层负责数据格式与加密压缩；应用层直接面向用户应用。
\begin{analysis}得分点：①ISO 提出、七层且层次名称与顺序正确（3 分，错一层扣 0.5 分）；②分层思想与优点（1 分）；③至少正确说明传输层、网络层、链路层三层功能（1 分）。常见失分：把 OSI 七层写成 TCP/IP 四层（网络接口层、网际层、运输层、应用层）；把“会话层、表示层”遗漏。\end{analysis}

\q \textbf{精准农业}\\
精准农业是以信息技术为支撑，以最小资源投入获得最大产出和最小环境负面影响为目标的现代农业生产管理体系。它的核心是按照田间小区乃至更小单元内土壤、作物、病虫害等条件的空间差异，实施定位、定时、定量的变量投入与管理，即“对症下药、按需供给”。其主要支撑技术是以遥感、地理信息系统、全球定位系统为核心的 3S 技术，以及传感器技术、变量作业机械和决策支持系统。实施过程一般包括农田信息采集、信息处理与差异分析、生成管理处方图、变量作业实施和效果评估五个环节。
\begin{analysis}得分点：①“按空间差异实施变量投入”“最小投入、最大产出、减少污染”各 2 分；②3S 技术与变量作业机械（1 分）。常见失分：把精准农业简单等同于“使用大型农机”或“用 GPS 开拖拉机”，未答出变量投入与处方图这一关键；与智慧农业完全不作区分。\end{analysis}

\newpage
\answerpart{二、简答题（每题 10 分，共 60 分）}

\q \textbf{信息系统的基本组成与主要功能}\\
（1）基本组成。信息系统是由人、硬件、软件、数据和处理规程五个要素有机结合而成的人机系统。硬件包括主机、存储设备、输入输出设备和网络设备；软件包括系统软件（操作系统、数据库管理系统）和应用软件（业务处理系统、决策支持系统）；数据是系统加工的对象和核心资源；人是系统的使用者和管理者；规程是保证系统正常运行的管理制度与操作规范。（4 分）

（2）主要功能。①信息的采集与输入，即把分散的原始数据收集并转化为系统可处理的形式；②信息的存储与组织，即对数据进行合理组织、维护和安全保护；③信息的加工处理，即对数据进行分类、排序、统计、检索、分析，得到对决策有用的信息；④信息的传输与输出，即通过网络把信息传递给使用者，并以报表、图表等形式呈现。（4 分）

（3）五个要素中，人是主导因素，数据是核心，硬件与软件是物质技术基础，规程是保障；只有五者协调配合，系统才能发挥整体效益。（2 分）
\begin{analysis}得分点：五要素齐全得 4 分（缺一扣 1 分）；四项功能齐全得 4 分；能指出人是主导、数据是核心、规程是保障得 2 分。常见失分：只答硬件和软件两项；把功能写成“输入、处理、输出”三项而不提存储与传输；把信息系统等同于计算机系统。\end{analysis}

\q \textbf{关系数据库中的视图及其作用}\\
（1）视图是从一个或多个基本表（或其他视图）导出的虚表，数据库中只存放视图的定义，不存放视图对应的数据；视图中的数据随基本表数据的变化而变化。（3 分）

（2）视图的主要作用：①简化操作，把复杂的多表连接和统计条件定义为一个视图，用户只需查询视图；②提供逻辑独立性，当基本表结构发生变化时，可通过修改视图定义使应用程序不必改动；③实现安全保护，只把视图中允许用户看到的行列暴露给用户，隐藏机密数据；④便于数据共享，同一基本表可以按不同用户的需要定义多种视图。（5 分）

（3）视图与基本表的区别：基本表是实际存储数据的表，视图是虚表，只保存定义不保存数据；对视图的更新最终要转换成对基本表的更新，且并非所有视图都可更新（含聚集函数、分组、多表连接等的视图通常不可更新）。（2 分）
\begin{analysis}得分点：视图定义“虚表、只存定义”得 3 分；作用答出简化操作、逻辑独立性、安全保护、数据共享四项得 5 分；区别得 2 分。常见失分：说视图一定物理保存一份数据；只答“方便查询”而漏掉安全性和逻辑独立性。\end{analysis}

\q \textbf{云计算 IaaS、PaaS、SaaS 三层服务方式}\\
（1）IaaS（基础设施即服务）：位于最底层，向用户提供计算、存储、网络等基础设施资源，用户通过虚拟化技术在云上获得虚拟机、云硬盘、虚拟网络，并自行安装操作系统和运行环境。特点是用户控制力最强、灵活性最高，但需要自行运维。典型应用是阿里云 ECS 弹性云服务器、亚马逊 AWS EC2。（3$\frac{1}{3}$ 分）

（2）PaaS（平台即服务）：位于中间层，在基础设施之上向用户提供应用开发与运行平台，包括操作系统、数据库、中间件、开发工具和运行环境。特点是用户只需关注应用逻辑，无须管理底层服务器和运行环境。典型应用是 Google App Engine、阿里云数据库和容器服务。（3$\frac{1}{3}$ 分）

（3）SaaS（软件即服务）：位于最上层，直接向用户提供可通过浏览器或客户端使用的应用软件，用户按需订购、按量付费，无须安装、部署和维护软件。典型应用是钉钉、企业微信、Office 365、Salesforce 等在线软件。（3$\frac{1}{3}$ 分）

三层的关系是自下而上的支撑关系：IaaS 提供“机房和机器”，PaaS 提供“平台和工具”，SaaS 提供“可以直接使用的软件”。
\begin{analysis}得分点：三层名称、层次位置、用户控制范围各 1 分，共 3 分；典型平台或应用举例 3 分（每层 1 分）；能说明三者自下而上的支撑关系 3 分。常见失分：把 PaaS 说成提供硬件；把 SaaS 举例为“租用服务器”；三层顺序颠倒。\end{analysis}

\q \textbf{Hadoop 的主要组件及其作用}\\
Hadoop 是 Apache 基金会推出的开源分布式计算平台，是大数据存储与批处理的代表性技术，主要由四个部分组成。（1 分）\\
（1）HDFS（Hadoop 分布式文件系统）：分布式存储组件，采用“一次写入、多次读取”模型，由 NameNode 保存文件系统的元数据并管理数据块的位置，由 DataNode 实际存储数据块，默认每块 128 MB 并保存多副本，从而提供高容错、高吞吐的存储能力。（3 分）\\
（2）MapReduce：分布式并行计算组件，把作业分解为 Map 和 Reduce 两个阶段，在数据所在结点上就近计算，实现大规模数据的并行处理。（2 分）\\
（3）YARN：集群资源管理与任务调度组件，由 ResourceManager 统一分配集群资源，由 NodeManager 管理单个结点的资源与容器，使多种计算框架可以共享同一集群。（2 分）\\
（4）Hadoop Common 与生态组件：Common 提供底层公共工具与接口；Hive 提供类 SQL 的数据仓库查询，HBase 提供列式实时读写，ZooKeeper 提供分布式协调，Sqoop 与 Flume 负责数据导入与采集。（2 分）
\begin{analysis}得分点：HDFS 及其 NameNode/DataNode 多副本机制得 3 分；MapReduce 得 2 分；YARN 得 2 分；Hive、HBase、ZooKeeper 等生态组件答出两个以上得 2 分；总体定位得 1 分。常见失分：把 MapReduce 与 Hadoop 等同，认为 Hadoop 就是 HDFS；把 YARN 说成数据库；把 NameNode 与 DataNode 的职责写反。\end{analysis}

\q \textbf{RFID 系统的组成、原理与优势}\\
（1）组成。RFID（射频识别）系统由电子标签、读写器和后端信息系统三部分组成。电子标签由芯片和天线构成，芯片内保存唯一标识及数据；读写器由天线、射频收发模块和控制单元构成，负责发射射频信号并接收标签返回的信息；后端系统（计算机与数据库）完成数据存储、处理与共享。按供电方式电子标签分为无源标签、半有源标签和有源标签。（3 分）\\
（2）工作原理。读写器通过天线发射特定频率的射频信号形成电磁场，进入场区的电子标签通过电磁耦合或电磁反向散射获得能量（无源标签）并被激活，把芯片中的数据调制后发射回读写器；读写器接收、解调并解码后送后端系统处理，从而完成非接触式的自动识别。（4 分）\\
（3）与条码技术相比的优势。可非接触、远距离读取；可穿透水、油、灰尘等遮蔽物读取；可同时批量识别多个标签；标签可反复读写、存储容量更大；寿命长、不易磨损，适应恶劣环境。因此在农产品溯源、仓储盘点、牲畜标识、冷链物流中具有明显优势。（3 分）
\begin{analysis}得分点：三部分组成 3 分；电磁耦合/反向散射、激活、调制回传的完整原理 4 分；优势答出非接触批量识别、可读写、抗污染三点得 3 分。常见失分：只答“电子标签和读卡器”而漏掉后端信息系统；把 RFID 说成需要人工对准扫描（那是条码）；原理只写“发射无线电波”而不写标签被激活并回传数据。\end{analysis}

\q \textbf{数据、信息、知识、智慧的概念与关系}\\
（1）概念。数据是对客观事物的性质、状态和相互关系进行记载的物理符号，是信息的载体，如数值、文字、声音、图像，本身是未加工的原料。信息是数据经过加工处理后、能够消除认识主体不确定性的内容，是数据的内容或诠释，如“土壤湿度 18\%”这一读数。知识是对大量信息进行归纳、提炼后形成的规律性认识，能回答“为什么”和“怎么办”，如“该作物在土壤含水量低于 20\% 时需灌溉”。智慧是综合运用知识和经验，结合具体情境进行判断、决策与创新的能力，如根据墒情、天气预报、作物长势和成本决定灌水时机与水量。（4 分）\\
（2）相互关系。四者是由低到高、逐级升华的认知层次：数据经过加工解释成为信息；信息经过归纳提炼成为知识；知识在具体情境中灵活运用上升为智慧。数据是基础，信息是加工结果，知识是规律化，智慧是创造性运用；层次越高，越依赖人的理解和判断，价值越大而获得难度也越高。（3 分）\\
（3）举例。农业生产环境监测中，传感器记录的“土壤湿度 18\%、气温 32$^\circ$C、光照 6 万 lx”是数据；处理后得到“土壤偏干、蒸发较快”是信息；把湿度、作物需水量与天气规律联系起来形成“含水量低于 20\% 即需灌溉”的判断是知识；进一步结合作物生育期、未来三天是否有雨、水电成本和产量目标，决定今天傍晚灌溉 15 mm，则是智慧。（3 分）
\begin{analysis}得分点：四个概念各 1 分（共 4 分）；逐级升华的关系表述 3 分；结合农业监测的完整举例 3 分。常见失分：把数据与信息混为一谈，只答“数据是数字、信息是消息”；举例时四个层次全写成同一内容，缺乏层级递进；漏掉智慧层。\end{analysis}

\newpage
\answerpart{三、论述题（每题 30 分，共 60 分）}

\q \textbf{我国农业信息化建设的主要问题与解决对策}\\
\textbf{一、破题与立意（约 3 分）。}农业信息化是把信息技术、信息资源与农业生产经营全过程深度融合，用信息化带动农业现代化的过程，是乡村振兴和农业强国建设的重要支撑。当前我国农业信息化已取得明显成效，但仍处于由起步向纵深推进的阶段，存在若干突出问题。

\textbf{二、存在的主要问题（约 12 分，每点 2 分，答出 6 点即可）。}\\
（1）农村信息基础设施薄弱且覆盖不均。部分偏远地区网络带宽不足、信号盲区多，田间地头的物联网基站、传感器供电与回传条件较差，难以支撑实时数据采集。\\
（2）数据资源分散、共享困难，形成“数据孤岛”。农业数据分属农业农村、气象、水利、市场监管等多个部门，标准不统一、接口不开放，存在重复采集与信息壁垒。\\
（3）农业信息标准体系不健全。传感器接口、数据格式、编码规则缺乏统一规范，不同厂商设备难以互联互通，影响系统集成与规模化推广。\\
（4）农业信息化人才严重短缺。农民整体信息素养和数字技能不高，基层农技人员既懂农业又懂信息技术的复合型人才匮乏，出现“设备建好了没人会用”的现象。\\
（5）资金投入不足且可持续性差。农业信息化项目公益性强、回报周期长，主要依赖财政投入，社会资本参与不足，建设完成后运维经费难以保障。\\
（6）农业信息服务体系不完善。信息发布不及时、针对性不强，农户“用不上、用不起、不会用”，信息服务与生产实际需求脱节。\\
（7）农业数据安全和隐私保护存在隐患。涉及耕地、产量、经营主体等敏感数据，安全防护和权属界定制度尚不健全。

\textbf{三、解决对策（约 12 分，每点 2 分，答出 6 点即可）。}\\
（1）加快农村信息基础设施建设。推进农村千兆光网、5G 和物联网覆盖，完善田间监测基站与电力保障，缩小城乡数字鸿沟，落实电信普遍服务。\\
（2）建设统一的农业大数据平台，打破数据壁垒。由主管部门牵头建立数据资源共享目录和交换标准，推动跨部门、跨层级数据汇聚共享，实现“一次采集、多方复用”。\\
（3）健全农业信息化标准体系。制定传感器接口、数据编码、平台接入和溯源标识等国家标准与行业标准，推行设备互联互通认证。\\
（4）加强人才培养与农民数字技能培训。依托高校开设农业信息工程相关专业，定向培养复合型人才；开展新型职业农民培训、科技特派员下乡，提高农户信息素养。\\
（5）创新投入机制。采取政府引导、企业主体、市场运作的模式，推广政府购买服务、PPP 和社会资本参与，建立信息化设施运维的长效经费保障机制。\\
（6）完善农业信息服务体系。发展“互联网＋农业”服务平台和农业专家系统，提供精准的气象、市场、病虫害与政策信息服务，实现由“信息推送”向“决策支持”升级。\\
（7）强化农业数据安全与治理。落实数据分类分级保护，加强身份认证、访问控制、加密传输和备份恢复，明确数据权属与使用规则。

\textbf{四、结合实际的深化分析（约 3 分）。}应坚持“需求导向、因地制宜、以用促建”的原则：东部发达地区可率先推进精准农业和智慧农场示范，中西部地区应优先解决网络覆盖与信息服务可及性问题；要防止重硬件轻应用、重建设轻运维的倾向，把农民的获得感和增收实效作为检验农业信息化成效的根本标准，最终实现以信息化驱动农业现代化、以数字化赋能乡村振兴。

\begin{analysis}得分点：问题部分答出 6 点以上且有具体内容得 12 分（每点 2 分，只写“基础薄弱、人才不足”等无展开的每点给 1 分）；对策部分答出 6 点以上且与问题对应得 12 分；立意与总体判断（农业信息化内涵、以信息化带动农业现代化）得 3 分；结合实际的原则性分析（因地制宜、以用促建、农民获得感）得 3 分。常见失分：只罗列问题不针对性地给对策，问题与对策不一一对应；对策空泛（“加大投入、加强领导”）而无具体措施；全篇不涉及我国农业农村实际，写成一般信息化概念题；只答“农村网络不好”，未涉及数据共享、标准、人才等深层问题。\end{analysis}

\q \textbf{人工智能技术及其在智慧农业中的应用}\\
\textbf{一、人工智能的基本概念与主要内容（约 8 分）。}人工智能（AI）是研究、开发用于模拟、延伸和扩展人的智能的理论、方法、技术及应用系统的技术科学，其目标是让机器具备感知、学习、推理、决策和交互的能力。主要技术内容包括：（1）知识表示与知识推理，用产生式规则、语义网络、本体等表示知识并进行逻辑推理，是专家系统的技术基础；（2）机器学习与深度学习，通过样本数据训练模型，实现分类、回归、聚类和预测，神经网络与深度学习是当前主流；（3）模式识别与计算机视觉，包括图像识别、目标检测、图像分割；（4）自然语言处理与语音识别；（5）智能控制与机器人技术；（6）智能 agent 与多智能体协同。（8 分）

\textbf{二、在智慧农业中的应用（约 14 分，每点 2 分，答出 7 点即可）。}\\
（1）病虫害识别与预警。利用卷积神经网络对叶片、果实的田间图像进行识别，秒级判断病虫害种类和危害程度，结合气象数据预测发生趋势，指导精准施药，减少农药用量。\\
（2）作物长势监测与产量预测。结合无人机和卫星遥感影像，通过深度学习模型反演叶面积指数、生物量，实现分区域长势评价和产量预估，为收储和销售提供依据。\\
（3）农业专家系统与智能决策。把农业专家的知识和经验建成知识库，通过推理机为农户提供品种选择、施肥配方、灌溉与防治方案，解决专家资源不足问题。\\
（4）智能灌溉与精准施肥。以土壤墒情、气象预报和作物需水规律为输入，用机器学习模型预测需水量，联动水肥一体化设备实现变量投入。\\
（5）畜禽养殖智能管理。通过个体识别、声音与行为分析、体温监测判断发情、疾病和应激状态，实现精准饲喂与环境自动调控。\\
（6）农产品分级与质量检测。基于机器视觉和近红外光谱对果蔬大小、色泽、损伤、糖度进行无损检测与自动分选。\\
（7）农机自动驾驶与农业机器人。融合北斗/GPS、视觉导航与智能控制，实现拖拉机自动导航、变量作业和采摘、除草机器人作业。\\
（8）农产品市场分析与农业信息服务。利用自然语言处理和预测模型分析价格行情、舆情与供需，为农户和政府部门提供决策支持。

\textbf{三、人工智能推动农业转型的作用分析（约 5 分）。}\\
人工智能使农业生产由依赖经验的“看天吃饭、凭经验种地”转向以数据为依据的“数据驱动、智能决策”：一是把分散的隐性经验固化为可复用的模型和知识库，解决农业专家稀缺问题；二是通过感知—分析—决策—执行的闭环，实现精准投入，提高水肥药利用率、降低成本并减少面源污染；三是提高对病虫害、气象灾害的预见性，增强抗风险能力；四是促进农业生产标准化、规模化和可追溯，提升农产品品质与品牌价值。

\textbf{四、面临的挑战与推进思路（约 3 分）。}\\
挑战主要有：农业数据获取成本高、样本少且标注质量参差，模型泛化能力不足；智能装备价格高、农村网络与电力条件受限；农民接受度和操作能力不足；农业数据涉及权属与隐私安全。推进思路是：建设农业大数据与标注样本库，推动产学研协同攻关适用模型；发展低成本、易操作的智能装备并做好农技培训与示范推广；健全农业数据安全与共享制度，坚持“人工智能＋农艺”融合，让人工智能真正服务于农业增效和农民增收。

\begin{analysis}得分点：人工智能概念与主要技术内容得 8 分（概念 3 分，技术内容答出 4 项以上得 5 分）；农业应用答出 6 个以上具体场景、每个有实质内容得 14 分（每点 2 分，只写“用于农业”无展开给 1 分）；对农业转型作用的分析得 5 分；挑战与推进思路得 3 分。常见失分：只罗列应用名称不作说明，成为名词堆砌；把人工智能等同于机器人或“计算机技术在农业的应用”；缺少分析和对策，写成简答题的放大版；完全没有结合我国农业实际（小农户为主、专家稀缺、面源污染等）。\end{analysis}

\end{document}
